% !TeX root = ../main.tex

\section{Обсуждение}

Бета-коэффициенты ОЛМ нельзя без дополнительных предпосылок считать прямым измерением силы нейронального ответа: они зависят от входного воздействия, гемодинамического преобразования, выбранной HRF и временной структуры исследования. Вычитание коэффициентов устраняет только общий аддитивный фон. Общий множитель сохраняется в масштабе контраста, а при нелинейном преобразовании ошибка зависит еще и от самих амплитуд. Это видно уже из зависимости $b(a)=a-\eta a^2$ с $\eta>0$: разность выходов равна $(a_A-a_B)[1-\eta(a_A+a_B)]=\Delta(1-2\eta m)$ и при неизменной разности амплитуд меняется с общим уровнем $m$ (ср. рис.~\ref{fig:f6}, \textit{c}). В динамической системе добавляется влияние предыстории: последовательность предъявления стимулов определяет состояние генератора и проекцию нелинейного остатка на коэффициенты условий. Поэтому конкретная последовательность событий важна не меньше средней частоты предъявления.

Точное линейное ядро и точная оценка конечной амплитуды решают разные задачи. $L_1$ описывает чувствительность модели около покоя, но не обязано минимизировать ошибки $\beta$ и контраста при конечном входе. Если проекция ошибки канонической формы противоположна нелинейному вкладу, оценка может оказаться ближе к заданному значению, однако такая компенсация зависит от амплитуд и последовательности предъявления стимулов. По той же причине хорошее воспроизведение BOLD-ряда не обеспечивает точность активационных оценок: расхождение сигнала и его проекция на выбранный контраст являются разными величинами.

Поправка второго порядка ограничена локальностью разложения и ошибкой оцененного входа. При росте амплитуды увеличивается роль высших порядков, поэтому правильное вычисление $Q_0$ не гарантирует правильную величину поправки. Если исходные оценки занижены, $\hat u_C$ описывает иную силу воздействия; если они содержат шум, квадратичное преобразование меняет и разброс, и среднее. Даже при нулевом среднем ошибки ее квадрат имеет положительное среднее; это известное смещение подстановочных оценок квадратичных функционалов~\cite{butucea2011}.

Две другие причины ухудшения точности связаны со шкалой наблюдения и несоответствием модели. При неизвестном выходном усилении номинальная и фактическая нелинейность масштабируются по-разному (подраздел <<Роль выходного усиления>>), что и проявилось при усилениях $0{,}8$ и $1{,}25$. Кроме того, номинальная модель может отличаться от генератора по параметрам или структуре. Сопоставление P-DCM и BW проверяет эту чувствительность, но не позволяет выбрать физиологически верное описание.

Результаты согласуются с представлением о BOLD как отклике нелинейной гемодинамической системы~\cite{friston2000,buxton2004} и с ролью нейрональных и сосудистых переходных процессов~\cite{havlicek2015}. Зависимость оценок от истории стимуляции отмечалась и в предшествующих работах~\cite{wager2005}. В настоящей работе количественно проверены канонические оценки и фиксированные поправки в проверенной модельной области. Использование собственного $L_1$ делает явным источник ошибки: видно, какая часть занижения остается после устранения рассогласования формы.

Систематическая составляющая ошибки важна и при усреднении индивидуальных оценок. Хотя групповое исследование не проводилось, из модели $\hat\Delta_i=\Delta_i+b_i+\varepsilon_i$ с независимым шумом нулевого среднего следуют два вывода. Смещение среднего равно $N^{-1}\sum_i b_i$, поэтому общий систематический сдвиг $b_i=b$ усреднением не уменьшается. Дисперсия шумовой части при одинаковой $\sigma^2$ составляет $\sigma^2/N$; если шумовые составляющие разных участников коррелированы, темп $1/\sqrt N$ не гарантирован. Усреднение индивидуально скорректированных оценок и поправка после усреднения BOLD являются разными операциями; эти аналитические следствия не заменяют групповую проверку.

Полученные данные не задают единого безопасного интервала между стимулами: для такого вывода нужно заранее определить допустимые ошибки $\beta$ и контраста, амплитуды, длительности, параметры генератора и правила агрегирования. Уменьшение ошибки при среднем SOA 8~с в рассмотренных сериях не устанавливает универсального порога 8~с для фМРТ. Уменьшение ошибок в одних условиях также не означает, что поправка безопасна в других; полезные группы, выделенные после теста, остаются описательными.

Генераторы имеют заданные параметры и не описывают пространственную структуру сосудистой сети, а общий линейный E/I-блок исключает нелинейную нейрональную адаптацию, которая может присутствовать в физиологических данных. Близость номинальной P-DCM к $C$ относится к конкретным параметрам. Переход к BW меняет несколько элементов сразу, поэтому различия результатов не идентифицируют единственный механизм.

Основные выводы получены для секундных событий и выбранных амплитуд. Блоковые дизайны и события другой длительности не исследовались, а отдельные серии используют разные последовательности предъявления стимулов и процедуры и не образуют единой проверенной поверхности. Поэтому численные границы нельзя переносить на другие дизайны без дополнительной проверки.

В шумовых сериях ковариация считалась известной, ее оценивание по данным не проверялось. Восемь повторов не дают точной оценки условных смещений, интервалы приближенные; мощность и частота ложных решений не исследованы. Независимой реализации внешним исследователем не было. Отрицательный итог FV на ее области не переносится автоматически на FW, а частные улучшения FW не отменяют ограничения FV.
